\documentclass[11pt]{article}

% Packages
\usepackage[margin=1in]{geometry}
\usepackage{amsmath}
\usepackage{listings}
\usepackage{xcolor}
\usepackage{hyperref}
\usepackage{natbib}

\lstdefinestyle{python}{
    language=Python,
    basicstyle=\ttfamily\small,
    keywordstyle=\color{blue},
    commentstyle=\color{gray},
    stringstyle=\color{red},
    numbers=left,
    numberstyle=\tiny\color{gray},
    stepnumber=1,
    numbersep=5pt,
    backgroundcolor=\color{white},
    showspaces=false,
    showstringspaces=false,
    showtabs=false,
    frame=single,
    tabsize=4,
    captionpos=b,
    breaklines=true,
    breakatwhitespace=false,
}

\lstdefinestyle{bash}{
    language=bash,
    basicstyle=\ttfamily\small,
    keywordstyle=\color{blue},
    commentstyle=\color{gray},
    stringstyle=\color{red},
    numbers=left,
    numberstyle=\tiny\color{gray},
    stepnumber=1,
    numbersep=5pt,
    backgroundcolor=\color{white},
    showspaces=false,
    showstringspaces=false,
    showtabs=false,
    frame=single,
    tabsize=4,
    captionpos=b,
    breaklines=true,
    breakatwhitespace=false,
}

\lstdefinestyle{sql}{
    language=SQL,
    basicstyle=\ttfamily\small,
    keywordstyle=\color{blue},
    commentstyle=\color{gray},
    stringstyle=\color{red},
    numbers=left,
    numberstyle=\tiny\color{gray},
    stepnumber=1,
    numbersep=5pt,
    backgroundcolor=\color{white},
    showspaces=false,
    showstringspaces=false,
    showtabs=false,
    frame=single,
    tabsize=4,
    captionpos=b,
    breaklines=true,
    breakatwhitespace=false,
}

\title{Engineering Runbook}
\author{Sakthyvell}

\begin{document}

\maketitle
\tableofcontents

\section{AWS for noobs}

\subsection{AWS Setup}



\section{Manual Testing as a SWE}
Manual testing is an art that I lack. As a software engineer, I have always focused on writing code and automating tests. However, I have come to realize the importance of manual testing in ensuring the quality of software. But the problem is, I suck at manual testing. I often miss edge cases and fail to think like an end-user. This has led to bugs and issues that could have been avoided with better manual testing. So this section is a growing list of test cases and scenarios that I need to remember while testing my code manually.

\subsection{Document Upload}
When testing document upload functionality, I need to consider the following scenarios:
\begin{itemize}
    \item Uploading valid documents of different formats (PDF, DOCX, TXT, etc.)
    \item Uploading invalid documents (corrupted files, unsupported formats)
    \item Uploading large files (test the maximum file size limit)
    \item Uploading multiple files at once
    \item Uploading files with special characters in the filename
    \item Uploading files with the same name (check for overwriting or versioning)
    \item Uploading files with different permissions (read-only, hidden files)
    \item Uploading files with different network conditions (slow connection, interrupted upload)
\end{itemize}

\subsection{API Endpoints}
When testing API endpoints, I need to consider the following scenarios:
\begin{itemize}
    \item Valid requests with correct parameters
    \item Invalid requests with missing or incorrect parameters
    \item Unauthorized access (invalid API keys, expired tokens)
    \item Rate limiting (exceeding the maximum number of requests)
    \item Testing different HTTP methods (GET, POST, PUT, DELETE)
    \item Testing response times and performance under load
    \item Testing error handling and response codes (400, 401, 403, 404, 500)
\end{itemize}



\bibliographystyle{plainnat}
\bibliography{refs}
\end{document}